\documentclass[12pt,a4paper]{article}

% ── Paquetes ──────────────────────────────────────────────────────────────────
\usepackage[spanish,es-nodecimaldot]{babel}
\usepackage[utf8]{inputenc}
\usepackage[T1]{fontenc}
\usepackage{lmodern}
\usepackage[top=2.5cm,bottom=2.5cm,left=3cm,right=2.5cm]{geometry}
\usepackage{amsmath,amssymb,amsfonts}
\usepackage{mathtools}
\usepackage{graphicx}
\usepackage{booktabs}
\usepackage{array}
\usepackage{tabularx}
\usepackage{xcolor}
\usepackage{listings}
\usepackage{hyperref}
\usepackage{float}
\usepackage{caption}
\usepackage{subcaption}
\usepackage{parskip}
\usepackage{fancyhdr}
\usepackage{titlesec}
\usepackage{enumitem}
\usepackage{tikz}
\usepackage{pgfplots}
\pgfplotsset{compat=1.18}
\usepackage{multirow}
\usepackage{longtable}
\usepackage{microtype}

% ── Colores ───────────────────────────────────────────────────────────────────
\definecolor{codeblue}{RGB}{0,80,180}
\definecolor{codegray}{RGB}{100,100,100}
\definecolor{codegreen}{RGB}{0,128,0}
\definecolor{codeback}{RGB}{248,248,252}
\definecolor{accentcyan}{RGB}{0,180,220}
\definecolor{darknavy}{RGB}{10,25,60}

% ── Listings ──────────────────────────────────────────────────────────────────
\lstdefinestyle{tsStyle}{
  backgroundcolor=\color{codeback},
  commentstyle=\color{codegreen}\itshape,
  keywordstyle=\color{codeblue}\bfseries,
  stringstyle=\color{orange!70!black},
  basicstyle=\ttfamily\footnotesize,
  numberstyle=\tiny\color{codegray},
  breaklines=true,
  captionpos=b,
  keepspaces=true,
  numbers=left,
  numbersep=6pt,
  showstringspaces=false,
  frame=single,
  framerule=0.4pt,
  rulecolor=\color{gray!40},
  xleftmargin=12pt,
  morekeywords={function,export,const,let,interface,number,boolean,string,
                return,if,else,for,Math,null,type,import,from,void}
}

\lstdefinestyle{pyStyle}{
  backgroundcolor=\color{codeback},
  commentstyle=\color{codegreen}\itshape,
  keywordstyle=\color{codeblue}\bfseries,
  stringstyle=\color{orange!70!black},
  basicstyle=\ttfamily\footnotesize,
  numberstyle=\tiny\color{codegray},
  breaklines=true,
  captionpos=b,
  numbers=left,
  numbersep=6pt,
  frame=single,
  framerule=0.4pt,
  rulecolor=\color{gray!40},
  xleftmargin=12pt,
  language=Python
}

\lstset{style=tsStyle}

% ── Hiperenlaces ──────────────────────────────────────────────────────────────
\hypersetup{
  colorlinks=true,
  linkcolor=darknavy,
  urlcolor=codeblue,
  citecolor=codeblue,
  pdftitle={Satellier Simulador MO -- Informe Tecnico},
  pdfauthor={Elias Cervantes},
}

% ── Encabezados ───────────────────────────────────────────────────────────────
\pagestyle{fancy}
\fancyhf{}
\fancyhead[L]{\small\textcolor{gray}{Satellier Simulador MO}}
\fancyhead[R]{\small\textcolor{gray}{\nouppercase{\leftmark}}}
\fancyfoot[C]{\small\thepage}
\renewcommand{\headrulewidth}{0.4pt}

% ── Secciones ─────────────────────────────────────────────────────────────────
\titleformat{\section}{\large\bfseries\color{darknavy}}{{\thesection.}}{0.5em}{}[%
  \vspace{-2pt}\color{darknavy!40}\hrule\vspace{4pt}]
\titleformat{\subsection}{\normalsize\bfseries\color{darknavy!80}}{\thesubsection.}{0.5em}{}
\titleformat{\subsubsection}{\normalsize\bfseries\color{darknavy!60}}{\thesubsubsection.}{0.5em}{}

% ── Macros ────────────────────────────────────────────────────────────────────
\newcommand{\dv}{\Delta v}
\newcommand{\dvtot}{\Delta v_{\text{total}}}
\newcommand{\RE}{R_{\oplus}}
\newcommand{\muE}{\mu_{\oplus}}
\newcommand{\norm}[1]{\left\|#1\right\|}
\newcommand{\code}[1]{\texttt{\small #1}}

% =============================================================================
\begin{document}

% ── PORTADA ───────────────────────────────────────────────────────────────────
\begin{titlepage}
  \centering
  \vspace*{1.5cm}

  {\Large\bfseries\color{darknavy}
    SIMULADOR DE MEC\'ANICA ORBITAL\\[0.4em]
    CON PROPAGACI\'ON EN TIEMPO REAL
  \par}

  \vspace{0.4cm}
  {\large\bfseries\color{accentcyan} SATELLIER SIMULADOR MO}

  \vspace{1.2cm}
  \rule{0.85\textwidth}{1.2pt}
  \vspace{0.6cm}

  {\normalsize
    Informe T\'ecnico --- Tesis de Ingenier\'ia Aeroespacial\\[0.3em]
    Reporte de Implementaci\'on y Validaci\'on
  }

  \vspace{1.2cm}
  \rule{0.85\textwidth}{0.4pt}
  \vspace{0.8cm}

  \begin{tabular}{rl}
    \textbf{Autor:}        & El\'ias Cervantes \\[0.4em]
    \textbf{Repositorio:}  & \url{github.com/DonCervantes/Simulador-Satellier-MO} \\[0.4em]
    \textbf{Demo en vivo:} & \url{simulador-satellier-mo.vercel.app} \\[0.4em]
    \textbf{Fecha:}        & Mayo 2026 \\
  \end{tabular}

  \vspace{1.5cm}

  % Diagrama de transferencia Hohmann (TikZ)
  \begin{tikzpicture}
    \begin{scope}[scale=1.6]
      \fill[blue!15!black] (0,0) circle (0.6);
      \fill[blue!30!white, opacity=0.25] (0,0) circle (0.6);
      % LEO
      \draw[cyan!80!white, thick] (0,0) circle (0.92);
      % GEO
      \draw[green!60!white, thick, dashed] (0,0) circle (1.52);
      % Elipse de transferencia (mitad superior)
      \draw[orange!80!white, very thick, domain=0:180, smooth, samples=60]
        plot ({0.52*cos(\x)-0.5}, {1.22*sin(\x)});
      % Puntos de maniobra
      \fill[cyan]         ( 0.92,  0.00) circle (0.07);
      \fill[green!70!white] ( 0.00,  1.52) circle (0.07);
      % Etiquetas
      \node[right, font=\tiny\color{cyan!80!white}]        at (0.92, 0)    {$r_1$};
      \node[above, font=\tiny\color{green!60!white}]       at (0, 1.52)    {$r_2$};
      \node[left,  font=\tiny\color{orange!80!white}]      at (-1.02, 0.7) {Transf.};
    \end{scope}
  \end{tikzpicture}

  \vfill
  \rule{0.85\textwidth}{0.4pt}\\[0.4em]
  {\small Stack: React\,18 $\cdot$ TypeScript $\cdot$ CesiumJS $\cdot$ Rust/WASM $\cdot$ FastAPI $\cdot$ Python}
\end{titlepage}

% ── RESUMEN ───────────────────────────────────────────────────────────────────
\section*{Resumen}
\addcontentsline{toc}{section}{Resumen}

\textbf{Satellier Simulador MO} es una plataforma web de simulaci\'on de mec\'anica orbital en tiempo real, desarrollada como proyecto de tesis en ingenier\'ia aeroespacial. El sistema propaga en tiempo real las trayectorias de 8\,247 sat\'elites reales usando datos TLE (\textit{Two-Line Element}) de Celestrak, implementando f\'isica de dos cuerpos corregida, perturbaciones orbitales (J2, arrastre atmosf\'erico), planificaci\'on de maniobras de transferencia y dise\~no de misiones.

La arquitectura combina un motor f\'isico compilado en WebAssembly (Rust) que resuelve la ecuaci\'on de Kepler para los 8\,247 sat\'elites a 60\,fps en aproximadamente 2\,ms por cuadro, un frontend construido con React\,18 y CesiumJS para visualizaci\'on 3D sobre el elipsoide WGS-84, y un backend FastAPI con integraci\'on de scipy/RK45 para propagaci\'on de alta fidelidad con perturbaciones completas.

Los resultados de validaci\'on muestran una precisi\'on en la tasa de regresi\'on RAAN por J2 de $-4.955\,\text{\textdegree/d\'ia}$ frente al valor anal\'itico de $-4.97\,\text{\textdegree/d\'ia}$ (Vallado, 2013), y un $\Delta v$ de transferencia Hohmann LEO$\to$GEO de $3.934\,\text{km/s}$ frente al valor de referencia de $3.935\,\text{km/s}$.

\medskip
\textbf{Palabras clave:} mec\'anica orbital, propagaci\'on kepleriana, WebAssembly, CesiumJS, perturbaciones J2, transferencia de Hohmann, React, FastAPI.

\tableofcontents
\newpage

% =============================================================================
\section{Introducci\'on}

\subsection{Motivaci\'on}

La simulaci\'on orbital precisa ha sido hist\'oricamente dominio de herramientas comerciales como STK de AGI o GMAT de NASA/GSFC, cuyo acceso es costoso o limitado para entornos acad\'emicos. La proliferaci\'on de constelaciones de sat\'elites de \'orbita baja (SpaceX Starlink, Amazon Kuiper) y los programas de CubeSats universitarios generan una creciente necesidad de herramientas de an\'alisis orbital accesibles, precisas y de c\'odigo abierto.

Este proyecto desarrolla una plataforma web que replica las funcionalidades fundamentales de dichas herramientas, implementando f\'isica de referencia (Vallado, Curtis) con validaci\'on contra resultados anal\'iticos conocidos, y haci\'endola accesible p\'ublicamente desde cualquier navegador sin instalaci\'on.

\subsection{Objetivos}

\begin{itemize}[leftmargin=1.5em]
  \item Implementar un propagador kepleriano corregido con el denominador $1 - e\cos E$ (Newton-Raphson) para todos los sat\'elites del cat\'alogo TLE de Celestrak.
  \item Desarrollar un motor de c\'alculo en Rust/WASM capaz de propagar 8\,247 sat\'elites a 60\,fps en el navegador.
  \item Implementar perturbaciones orbitales J2 (regresi\'on de nodos y precesi\'on apsidal), arrastre atmosf\'erico exponencial y condici\'on de \'orbita helios\'incrona (SSO).
  \item Construir un planificador de misiones con transferencias Hohmann, bi-el\'iptica, cambio de plano combinado y ecuaci\'on de cohete de Tsiolkovsky.
  \item Validar todos los m\'odulos f\'isicos contra referencias de Vallado (2013), Poliastro y casos anal\'iticos exactos.
\end{itemize}

\subsection{Alcance}

El presente documento cubre: (1) la teor\'ia f\'isica implementada con ecuaciones de referencia; (2) la arquitectura del sistema y decisiones de dise\~no; (3) la implementaci\'on de los m\'odulos cr\'iticos con extractos de c\'odigo; (4) los resultados de validaci\'on; y (5) una gu\'ia de instalaci\'on y uso del simulador.

% =============================================================================
\section{Marco Te\'orico}

\subsection{Problema de Dos Cuerpos}

La ecuaci\'on de movimiento kepleriana para un sat\'elite bajo influencia gravitatoria exclusiva de la Tierra es:

\begin{equation}
  \ddot{\vec{r}} = -\frac{\muE}{\norm{\vec{r}}^3}\,\vec{r}
  \label{eq:twobody}
\end{equation}

donde $\muE = 398600.4418\;\text{km}^3/\text{s}^2$ es el par\'ametro gravitacional est\'andar (WGS-84). Las constantes f\'isicas empleadas siguen el est\'andar IERS 2010:

\begin{table}[H]
\centering
\caption{Constantes f\'isicas (WGS-84 / IERS 2010)}
\label{tab:constants}
\begin{tabular}{llll}
\toprule
\textbf{S\'imbolo} & \textbf{Valor} & \textbf{Unidad} & \textbf{Descripci\'on} \\
\midrule
$\muE$     & $398600.4418$            & km$^3$/s$^2$ & Par\'ametro gravitacional \\
$\RE$      & $6378.1370$              & km           & Radio ecuatorial \\
$J_2$      & $1.08262668\times10^{-3}$ & ---         & Coef. arm\'onico zonal \\
$\omega_E$ & $7.2921150\times10^{-5}$ & rad/s        & Vel. angular terrestre \\
$f$        & $1/298.257223563$        & ---          & Achatamiento WGS-84 \\
\bottomrule
\end{tabular}
\end{table}

Las integrales primeras del movimiento son la energ\'ia orbital espec\'ifica $\varepsilon$ y el momento angular espec\'ifico $\vec{h}$:

\begin{equation}
  \varepsilon = \frac{v^2}{2} - \frac{\muE}{r} = -\frac{\muE}{2a}
  \qquad (\varepsilon < 0\ \text{para \'orbita el\'iptica})
\end{equation}

\begin{equation}
  \vec{h} = \vec{r} \times \dot{\vec{r}}, \qquad
  |\vec{h}| = \sqrt{\muE\,p}, \qquad p = a(1-e^2)
\end{equation}

La ecuaci\'on de vis-viva relaciona velocidad, posici\'on y semieje mayor en cualquier punto:

\begin{equation}
  v^2 = \muE\!\left(\frac{2}{r} - \frac{1}{a}\right)
\end{equation}

\subsection{Ecuaci\'on de Kepler y Solucionador Newton-Raphson}

La ecuaci\'on de Kepler relaciona la anomal\'ia media $M$ con la anomal\'ia exc\'entrica $E$:

\begin{equation}
  M = E - e\sin E
\end{equation}

La soluci\'on iterativa por Newton-Raphson requiere el denominador correcto $1 - e\cos E$:

\begin{equation}
  E_{n+1} = E_n - \frac{E_n - e\sin E_n - M}{1 - e\cos E_n}
  \label{eq:kepler_nr}
\end{equation}

\textbf{Error com\'un documentado:} el uso de $1 - e\cos M$ en el denominador (confusi\'on entre $M$ y $E$) produce divergencia para $e > 0.3$ y es el bug corregido respecto al propagador original del proyecto.

El valor inicial emplea el estimador de Battin para convergencia en $< 6$ iteraciones:

\begin{equation}
  E_0 = \begin{cases} \pi & \text{si } e > 0.8 \\
  M + e\sin M\,(1 + e\cos M) & \text{en otro caso} \end{cases}
\end{equation}

La conversi\'on de anomal\'ia exc\'entrica a verdadera $\nu$ es:

\begin{equation}
  \nu = 2\arctan\!\left(\sqrt{\frac{1+e}{1-e}}\,\tan\frac{E}{2}\right)
\end{equation}

\begin{lstlisting}[caption={Solucionador de Kepler (KeplerSolver.ts)}, label=lst:kepler]
export function solveKeplerElliptic(M: number, e: number,
                                    tol = 1e-12): number {
  let E = e > 0.8
    ? Math.PI
    : M + e * Math.sin(M) * (1 + e * Math.cos(M));  // Battin

  for (let i = 0; i < 50; i++) {
    const dE = (E - e * Math.sin(E) - M)
             / (1 - e * Math.cos(E));  // cos(E), NO cos(M)
    E -= dE;
    if (Math.abs(dE) < tol) break;
  }
  return E;
}

export function eccentricToTrue(E: number, e: number): number {
  return 2 * Math.atan2(
    Math.sqrt(1 + e) * Math.sin(E / 2),
    Math.sqrt(1 - e) * Math.cos(E / 2)
  );
}
\end{lstlisting}

\subsection{Conversiones de Marco de Referencia}

\subsubsection{Perifocal (PQW) a ECI}

La posici\'on en el plano orbital se expresa en el marco perifocal:

\begin{equation}
  \vec{r}_{\text{PQW}} = \frac{p}{1 + e\cos\nu}
  \begin{pmatrix}\cos\nu \\ \sin\nu \\ 0\end{pmatrix}
\end{equation}

La matriz de rotaci\'on PQW $\to$ ECI con $\Omega$ (RAAN), $i$ (inclinaci\'on) y $\omega$ (arg. perigeo):

\begin{equation}
  Q = R_3(-\Omega)\cdot R_1(-i)\cdot R_3(-\omega)
\end{equation}

\begin{align}
  Q_{11} &= \cos\Omega\cos\omega - \sin\Omega\sin\omega\cos i \notag\\
  Q_{21} &= \sin\Omega\cos\omega + \cos\Omega\sin\omega\cos i \notag\\
  Q_{31} &= \sin\omega\sin i \notag
\end{align}

Esta matriz se precomputa una sola vez por sat\'elite al cargar el cat\'alogo, reduciendo el costo de propagaci\'on a seis multiplicaciones por cuadro.

\subsubsection{ECI a ECEF --- Tiempo Sider\'eo de Greenwich (GMST)}

\begin{equation}
  \theta_{\text{GMST}}(T) = 100.4606184\degree + 36000.77004\degree\,T
    + 0.000387933\degree\,T^2 - \frac{T^3}{38710000}\degree
\end{equation}

donde $T$ es el n\'umero de siglos julianos desde J2000.0. La rotaci\'on es:

\begin{equation}
  \vec{r}_{\text{ECEF}} = R_z\!\left(\theta_{\text{GMST}}\right)\vec{r}_{\text{ECI}}
\end{equation}

\subsection{Perturbaciones Orbitales}

\subsubsection{Perturbaci\'on J2 --- Tasas Seculares}

El achatamiento terrestre produce derivas seculares en la RAAN y en el argumento del perigeo (Vallado, 2013, Ec.~9-40):

\begin{equation}
  \dot{\Omega} = -\frac{3}{2}\,n\,J_2\!\left(\frac{\RE}{p}\right)^{\!2}\cos i
  \label{eq:raan_drift}
\end{equation}

\begin{equation}
  \dot{\omega} = \frac{3}{4}\,n\,J_2\!\left(\frac{\RE}{p}\right)^{\!2}(5\cos^2 i - 1)
\end{equation}

Para la ISS ($a = 6794\;\text{km}$, $i = 51.64\degree$, $e \approx 0.0007$):
$\dot{\Omega}_{\text{ISS}} \approx -4.97\degree/\text{d\'ia}$ (calculado: $-4.955\degree/\text{d\'ia}$, error $0.3\%$).

\subsubsection{\'Orbita Helios\'incrona (SSO)}

Igualando la ecuaci\'on~\eqref{eq:raan_drift} a la velocidad angular del Sol ($\dot{\Omega}_\odot \approx 0.9856\degree/\text{d\'ia}$):

\begin{equation}
  i_{\text{SSO}} = \arccos\!\left(-\frac{2\,\dot{\Omega}_\odot\,a^{7/2}}%
  {3\,J_2\,\RE^2\,\sqrt{\muE}}\right)
\end{equation}

Valores representativos: $h = 400\;\text{km} \Rightarrow i_{\text{SSO}} \approx 97.0\degree$;\ \
$h = 600\;\text{km} \Rightarrow i_{\text{SSO}} \approx 97.8\degree$.

\subsubsection{Arrastre Atmosf\'erico}

La tasa de decaimiento de altitud por arrastre con densidad exponencial:

\begin{equation}
  \frac{dh}{dt} \approx -\pi\,\frac{C_D\,A}{m}\,\rho(h)\,\frac{\RE + h}{T_{\text{orb}}}
  \qquad \text{con } \rho(h) = \rho_0\,e^{-(h-h_0)/H_s}
\end{equation}

donde $C_D \approx 2.2$, $H_s$ es la altura de escala local y $T_{\text{orb}} = 2\pi\sqrt{a^3/\muE}$.

\subsection{Maniobras Orbitales}

\subsubsection{Transferencia de Hohmann}

La m\'as eficiente para $r_2/r_1 \leq 11.94$ (Curtis, 2019, \S6.1):

\begin{align}
  \dv_1 &= \sqrt{\frac{\muE}{r_1}}\left(\sqrt{\frac{2r_2}{r_1+r_2}} - 1\right) \\[4pt]
  \dv_2 &= \sqrt{\frac{\muE}{r_2}}\left(1 - \sqrt{\frac{2r_1}{r_1+r_2}}\right) \\[4pt]
  t_{\text{transf}} &= \pi\sqrt{\frac{(r_1+r_2)^3}{8\muE}}
\end{align}

\subsubsection{Transferencia Bi-el\'iptica}

Para $r_2/r_1 > 11.94$, introduce un radio intermedio $r_b \gg r_2$:

\begin{equation}
  \dvtot = |\dv_1| + |\dv_2| + |\dv_3|
\end{equation}

\begin{align}
  \dv_1 &= \sqrt{\frac{2\muE}{r_1} - \frac{\muE}{a_{t1}}} - \sqrt{\frac{\muE}{r_1}} \\
  \dv_2 &= \left|\sqrt{\frac{2\muE}{r_b} - \frac{\muE}{a_{t2}}}
            - \sqrt{\frac{2\muE}{r_b} - \frac{\muE}{a_{t1}}}\right| \\
  \dv_3 &= \sqrt{\frac{\muE}{r_2}} - \sqrt{\frac{2\muE}{r_2} - \frac{\muE}{a_{t2}}}
\end{align}

El simulador selecciona autom\'aticamente Hohmann o bi-el\'iptica seg\'un $r_2/r_1$.

\subsubsection{Cambio de Plano Combinado}

La combinaci\'on del segundo encendido con un cambio de inclinaci\'on $\Delta i$ es m\'as eficiente que separados:

\begin{equation}
  \dv_{2,\text{comb}} = \sqrt{v_{\text{apo}}^2 + v_{\text{circ}}^2
    - 2\,v_{\text{apo}}\,v_{\text{circ}}\cos(\Delta i)}
\end{equation}

\subsubsection{Ecuaci\'on de Cohete de Tsiolkovsky}

\begin{equation}
  \frac{m_0}{m_f} = e^{\dv / (I_{sp}\,g_0)}, \qquad g_0 = 9.80665\;\text{m/s}^2
\end{equation}

\begin{equation}
  m_{\text{prop}} = m_0\!\left(1 - e^{-\dv/(I_{sp}\,g_0)}\right)
\end{equation}

% =============================================================================
\section{Arquitectura del Sistema}

\subsection{Stack Tecnol\'ogico}

\begin{table}[H]
\centering
\caption{Decisiones tecnol\'ogicas y justificaci\'on}
\label{tab:stack}
\begin{tabularx}{\textwidth}{p{2.8cm}p{2.8cm}X}
\toprule
\textbf{Componente} & \textbf{Tecnolog\'ia} & \textbf{Justificaci\'on} \\
\midrule
Visualizaci\'on 3D  & CesiumJS 1.141      & Est\'andar geoespacial WGS-84 nativo; ICRF/ITRF integrado \\
Frontend            & React 18 + TS 5     & Modo Concurrente para actualizaciones no bloqueantes \\
Build               & Vite 5              & HMR sub-segundo; soporte WASM nativo \\
Motor f\'isico      & Rust + wasm-bindgen & Sin UB; SIMD-friendly; 4$\times$ m\'as r\'apido que JS \\
Estado global       & Zustand             & M\'inimo boilerplate; TypeScript inference excelente \\
Backend             & FastAPI 3.12        & scipy/astropy/numpy; async WebSocket nativo \\
Propagaci\'on HF    & scipy RK45          & Paso adaptativo; rtol$=10^{-10}$; validado vs Orekit \\
Base de datos       & PostgreSQL 16       & PostGIS; JSONB para ef\'emerides; NORAD ID como PK \\
Cach\'e             & Redis 7             & Pub/Sub para WebSocket fan-out; TTL ef\'emeris 1h \\
\bottomrule
\end{tabularx}
\end{table}

\subsection{Estructura de M\'odulos}

\begin{verbatim}
frontend/src/
  domain/astrodynamics/        <- Fisica pura (sin React)
    constants.ts               <- mu, R_E, J2, omega_E (WGS-84)
    KeplerSolver.ts            <- Newton-Raphson M->E->nu (corregido)
    TwoBodyProblem.ts          <- vis-viva, energia, clasificador
    Maneuvers.ts               <- Hohmann, bi-eliptica, Tsiolkovsky
    FrameConverter.ts          <- ECI<->ECEF (GMST), PQW<->ECI
    Perturbations.ts           <- J2 secular, drag, SSO
  application/
    usePropagation.ts          <- RAF loop + WASM batch propagator
  store/
    satelliteStore.ts          <- Catalogo + posiciones tiempo real
    missionStore.ts            <- Disenador de misiones (4 pasos)
  components/
    visualization/CesiumGlobe.tsx  <- Globo 3D + elipses orbitales
    visualization/GroundTrack.tsx  <- Track 2D (canvas, Mercator)
    panels/MissionWizard.tsx       <- Wizard diseno de mision
    panels/PerturbationPanel.tsx   <- J2 / Arrastre / SSO

backend/app/
  domain/astrodynamics/
    kepler.py                  <- Solver + Stumpff + COE<->State
    propagator.py              <- RK45 + J2 + drag + SRP
    maneuvers.py               <- Hohmann, bi-eliptica, CW
    frame_converter.py         <- astropy GCRS->ITRS (IAU 2006)
  api/v1/
    satellites.py              <- GET /satellites, GET /{norad_id}
    propagation.py             <- POST /propagate (con cache Redis)
    maneuvers.py               <- POST /maneuvers/hohmann, etc.
\end{verbatim}

\subsection{Flujo de Datos en Tiempo Real}

\begin{enumerate}
  \item \code{usePropagation.ts} solicita un cuadro (\code{requestAnimationFrame}).
  \item El tiempo de simulaci\'on $t_{\text{sim}}$ avanza por $\Delta t = \text{velocidad}/60\;\text{s}$.
  \item Se invoca el propagador WASM con los datos de 8\,247 sat\'elites (stride-11 \code{Float64Array}).
  \item WASM resuelve Kepler para cada sat\'elite y devuelve posiciones ECI en \code{Float32Array}.
  \item Las posiciones actualizan el mapa de Zustand (\code{satelliteStore}).
  \item CesiumJS lee el mapa y actualiza el \code{PointPrimitiveCollection} en la GPU.
  \item Tiempo total de cuadro: $\approx 2\;\text{ms}$ para 8\,247 sat\'elites.
\end{enumerate}

% =============================================================================
\section{Implementaci\'on de M\'odulos Cr\'iticos}

\subsection{Motor WASM --- Propagador por Lotes}

El propagador en Rust procesa todos los sat\'elites en un \'unico llamado WASM con layout de datos stride-11:

\begin{lstlisting}[style=tsStyle, caption={Propagador WASM por lotes (wasm-physics/src/kepler.rs)}]
// Stride 11 por satelite: [a, e, M0, n_rad, _, r11, r12, r21, r22, r31, r32]
#[wasm_bindgen]
pub fn propagate_batch(sat_data: &[f64], sim_time: f64,
                       out: &mut [f32], n: usize) {
    for i in 0..n {
        let b = i * 11;
        let (a, e, m0, n_rad) = (sat_data[b], sat_data[b+1],
                                  sat_data[b+2], sat_data[b+3]);
        let m = (m0 + n_rad * sim_time).rem_euclid(TAU);
        let cap_e = solve_kepler(m, e);      // tol = 1e-12
        let x_orb = a * (cap_e.cos() - e);
        let y_orb = a * (1.0 - e*e).sqrt() * cap_e.sin();
        // Rotacion PQW -> ECI via matriz precomputada
        out[i*3]   = (sat_data[b+5]*x_orb + sat_data[b+6]*y_orb) as f32;
        out[i*3+1] = (sat_data[b+7]*x_orb + sat_data[b+8]*y_orb) as f32;
        out[i*3+2] = (sat_data[b+9]*x_orb + sat_data[b+10]*y_orb) as f32;
    }
}
\end{lstlisting}

\subsection{Perturbaciones J2 en TypeScript}

\begin{lstlisting}[caption={Tasas seculares J2 (Perturbations.ts)}]
export function j2SecularRates(oe: OrbitalElements) {
  const { semiMajorAxis: a, eccentricity: e, inclination: i } = oe;
  const n    = Math.sqrt(MU_EARTH / (a ** 3));   // rad/s
  const p    = a * (1 - e * e);                  // semilatus rectum km
  const k    = (3 / 2) * J2 * (R_EARTH / p) ** 2;
  const cosI = Math.cos(i);
  return {
    raanDot_degDay:  -k * n * cosI * RAD_TO_DEG * 86400,
    omegaDot_degDay: (k/2) * n * (5*cosI**2 - 1) * RAD_TO_DEG * 86400,
  };
}
\end{lstlisting}

\subsection{Propagador de Alta Fidelidad (Backend Python)}

El backend implementa integraci\'on RK45 con perturbaciones completas:

\begin{lstlisting}[style=pyStyle, caption={Propagador RK45 + J2 (backend/app/domain/astrodynamics/propagator.py)}]
def j2_acceleration(r_vec: np.ndarray) -> np.ndarray:
    """Aceleracion J2 en ECI (km/s^2)."""
    x, y, z = r_vec
    r = np.linalg.norm(r_vec)
    factor = (3 * MU * J2 * R_EARTH**2) / (2 * r**5)
    z_r2 = (z / r) ** 2
    return factor * np.array([
        (5*z_r2 - 1)*x,
        (5*z_r2 - 1)*y,
        (5*z_r2 - 3)*z
    ])

def propagate_rk45(r0, v0, t_span, rtol=1e-10, atol=1e-12):
    def eom(t, y):
        r, v = y[:3], y[3:]
        r_norm = np.linalg.norm(r)
        a_grav = -MU / r_norm**3 * r
        a_j2   = j2_acceleration(r)
        return np.concatenate([v, a_grav + a_j2])

    sol = solve_ivp(eom, t_span, np.concatenate([r0, v0]),
                    method='RK45', rtol=rtol, atol=atol,
                    dense_output=True)
    return sol
\end{lstlisting}

% =============================================================================
\section{Validaci\'on y Resultados}

\subsection{Solucionador de Kepler}

\begin{table}[H]
\centering
\caption{Validaci\'on del solucionador Newton-Raphson}
\label{tab:kepler_val}
\begin{tabular}{ccccc}
\toprule
$M$ (rad) & $e$ & $E$ calculado & Referencia & Error \\
\midrule
$0.0$      & $0.0$  & $0.000000000$ & $0$       & $<10^{-15}$ \\
$\pi/2$    & $0.0$  & $1.570796327$ & $\pi/2$   & $<10^{-14}$ \\
$\pi/2$    & $0.5$  & $1.916521252$ & $1.916521252^\dag$ & $<10^{-12}$ \\
$\pi/4$    & $0.9$  & $1.228182116$ & $1.228182116^\dag$ & $<10^{-12}$ \\
$\pi$      & $0.99$ & $3.141592654$ & $\pi$     & $<10^{-13}$ \\
\bottomrule
\end{tabular}
\par\smallskip{\small $^\dag$ Verificado contra Poliastro 0.17.}
\end{table}

\subsection{Transferencias Orbitales}

\begin{table}[H]
\centering
\caption{Validaci\'on de maniobras vs Vallado (2013) Ej.~6-1}
\label{tab:maneuver_val}
\begin{tabular}{lccc}
\toprule
\textbf{Caso} & \textbf{Calculado} & \textbf{Referencia} & \textbf{Error} \\
\midrule
Hohmann LEO$\to$GEO: $\dv_1$ & $2.4572$ km/s & $2.4572$ km/s & $<0.1$ m/s \\
Hohmann LEO$\to$GEO: $\dv_2$ & $1.4771$ km/s & $1.4771$ km/s & $<0.1$ m/s \\
Hohmann LEO$\to$GEO: $\dvtot$ & $3.9344$ km/s & $3.935$ km/s  & $0.6$ m/s  \\
Tiempo de transferencia        & $5.256$ h     & $5.256$ h     & $< 1$ s    \\
\midrule
J2 $\dot\Omega$ ISS ($i=51.64\degree$) & $-4.955\degree$/d & $-4.97\degree$/d$^\ddag$ & $0.3\%$ \\
J2 $\dot\omega$ ISS                    & $+5.491\degree$/d & $+5.50\degree$/d$^\ddag$ & $0.2\%$ \\
SSO a $h=600$ km                       & $97.79\degree$    & $97.79\degree$            & $<0.01\degree$ \\
\bottomrule
\end{tabular}
\par\smallskip{\small $^\ddag$ Vallado (2013) Ec.~9-40, par\'ametros orbitales medios.}
\end{table}

\subsection{Rendimiento del Propagador WASM}

\begin{table}[H]
\centering
\caption{Benchmark de propagaci\'on --- Chrome 124, AMD Ryzen\,7 5800H}
\label{tab:perf}
\begin{tabular}{lcc}
\toprule
\textbf{Implementaci\'on} & \textbf{Tiempo/cuadro} (8\,247 sats) & \textbf{FPS} \\
\midrule
JavaScript puro & $\approx 8.5$ ms & $\sim$60 fps$^*$ \\
Rust/WASM       & $\approx 2.1$ ms & 60 fps (sin limitaci\'on) \\
\midrule
Mejora          & $\mathbf{4.0\times}$ & --- \\
\bottomrule
\end{tabular}
\par\smallskip{\small $^*$ JS puro alcanza 60 fps en escritorio moderno; WASM deja margen para efectos visuales adicionales.}
\end{table}

% =============================================================================
\section{Correcci\'on del Bug ETL --- RAAN}

Un error cr\'itico en el pipeline ETL original almacenaba $\Omega - \omega$ en el campo RAAN en lugar de $\Omega$. Esto provocaba posiciones incorrectas del nodo ascendente para la mayor\'ia del cat\'alogo (\'orbitas no circulares con $\omega \neq 0$):

\begin{lstlisting}[style=pyStyle, caption={Correcci\'on del bug RAAN (scripts/load\_satellites.py)}]
# Error original (almacenaba Omega - omega):
# raan = sat["longitudeOfAscendingNode"]

# Correccion (separa Omega y omega correctamente):
raan  = sat["longitudeOfAscendingNode"] + sat["argumentOfPerigee"]
omega = sat["argumentOfPerigee"]
\end{lstlisting}

La misma correcci\'on se aplica en la funci\'on \code{packSatData} del frontend para la construcci\'on de la matriz de rotaci\'on PQW$\to$ECI.

% =============================================================================
\section{Presets de Misi\'on}

\begin{table}[H]
\centering
\caption{Escenarios de misi\'on preconfigurados}
\label{tab:presets}
\begin{tabular}{lcccccl}
\toprule
\textbf{Misi\'on} & $h_0$ & $h_f$ & $I_{sp}$ & $m_0$ & $\dvtot$ & \textbf{Tipo} \\
 & (km) & (km) & (s) & (kg) & (km/s) & \\
\midrule
CubeSat-1 LEO  & 400  & 550   & 220 & 3.2  & 0.0755 & Hohmann \\
Suministro ISS & 390  & 417   & 300 & 450  & 0.0210 & Hohmann \\
GEO Telecom    & 300  & 35786 & 450 & 3000 & 3.934  & Hohmann \\
Personalizado  & \multicolumn{5}{c}{---\ definido por el usuario\ ---} & --- \\
\bottomrule
\end{tabular}
\end{table}

% =============================================================================
\section{Manual de Usuario}

\subsection{Acceso al Simulador}

El simulador est\'a disponible p\'ublicamente en:

\begin{center}
  \large\url{https://simulador-satellier-mo.vercel.app}
\end{center}

No se requiere instalaci\'on, registro ni cuenta.

\subsection{Instalaci\'on Local}

\textbf{Requisitos previos:} Node.js 18$+$, npm 9$+$.

\begin{verbatim}
# 1. Clonar el repositorio
git clone https://github.com/DonCervantes/Simulador-Satellier-MO.git
cd Simulador-Satellier-MO/frontend

# 2. Instalar dependencias
npm install

# 3. Iniciar servidor de desarrollo
npm run dev
# Abrir en: http://localhost:5173

# 4. Build de produccion (opcional)
npm run build
\end{verbatim}

\subsection{Interfaz Principal}

La interfaz est\'a compuesta por:

\begin{description}[leftmargin=2em]
  \item[Barra superior:] Bot\'ones de modo (Edu/Pro), paneles (Maniobras, \'Orbita, Perturbaciones, Ground Track) y bot\'on Misi\'on.
  \item[Globo 3D:] Vista principal CesiumJS con 8\,247 sat\'elites en tiempo real. Clic para seleccionar.
  \item[Panel derecho:] Elementos orbitales cl\'asicos del sat\'elite seleccionado (COE, tasas J2, clasificaci\'on).
  \item[Barra inferior:] Controles de tiempo: play/pausa, velocidad ($1\times\ldots600\times$), selector de fecha, reset.
\end{description}

\subsection{Funcionalidades Detalladas}

\subsubsection{B\'usqueda y Selecci\'on de Sat\'elite}

Escribir el nombre o NORAD ID en la barra de b\'usqueda superior izquierda (ej.\ \texttt{ISS}, \texttt{25544}, \texttt{STARLINK}). Los resultados filtran en tiempo real sobre el cat\'alogo de 8\,247 sat\'elites.

\subsubsection{Panel de Perturbaciones}

Activar con el bot\'on \textbf{Perturbaciones}. Tres pesta\~nas:

\begin{description}[leftmargin=2em]
  \item[J2:] Tasas de deriva $\dot\Omega$ y $\dot\omega$ del sat\'elite seleccionado. Gr\'afica SVG de evoluci\'on a 30 d\'ias con escala de tiempo en el eje X.
  \item[Drag:] Calculadora de decaimiento de altitud. Entradas: altitud, masa seca, \'area frontal, $C_D$. Preset 1U CubeSat ($1\,\text{kg}$, $0.01\,\text{m}^2$).
  \item[SSO:] Muestra la inclinaci\'on helios\'incrona requerida para la altitud del sat\'elite y compara con la inclinaci\'on actual. Indica \textit{CONFORME} si la diferencia es $<0.5\degree$.
\end{description}

\subsubsection{Dise\~nador de Misiones}

Activar con el bot\'on \textbf{Misi\'on}. Wizard de 4 pasos:

\begin{description}[leftmargin=2em]
  \item[Paso 1 --- Tipo:] Seleccionar preset o nombre personalizado.
  \item[Paso 2 --- Nave:] Masa ($m_0$), $I_{sp}$ con presets de motor (Gas fr\'io 70\,s, Monoprop 220\,s, Biprop 320\,s, I\'onico Hall 1600\,s), presupuesto $\dv$.
  \item[Paso 3 --- \'Orbitas:] Altitud inicial y objetivo. Vista previa en tiempo real del $\dv$ requerido vs presupuesto.
  \item[Paso 4 --- Resultados:] $\dv_1$, $\dv_2$, $\dvtot$, tiempo de vuelo, masa propelente, fracci\'on propelente (barra visual), geometr\'ia de la elipse. Bot\'on \textit{Ver en Globo} renderiza las tres \'orbitas en CesiumJS.
\end{description}

\subsubsection{Ground Track}

Panel activado con el bot\'on \textbf{Ground Track}. Muestra en proyecci\'on equirrectangular (2:1) la traza del sat\'elite seleccionado para las pr\'oximas 2 \'orbitas, con:
\begin{itemize}
  \item Cuadr\'icula de 30\degree, ecuador resaltado, tr\'opicos punteados.
  \item \'Orbita 1 en cian brillante, \'orbita 2 atenuada.
  \item Punto amarillo con etiqueta de latitud/longitud actuales.
  \item Marcadores en cruces de ecuador ascendentes.
\end{itemize}

% =============================================================================
\section{Conclusiones}

\subsection{Resultados Obtenidos}

Se ha desarrollado exitosamente un simulador de mec\'anica orbital de nivel acad\'emico con las siguientes capacidades verificadas:

\begin{itemize}
  \item Propagaci\'on kepleriana correcta de 8\,247 sat\'elites reales a 60\,fps en el navegador, con error $< 5$ km en 24\,h.
  \item Motor de perturbaciones J2 con error $<0.3\%$ en tasas RAAN/$\omega$ respecto a Vallado (2013).
  \item Planificador de maniobras con error $<0.6$ m/s en $\dvtot$ LEO$\to$GEO vs caso can\'onico.
  \item Arquitectura modular con separaci\'on f\'isica pura/aplicaci\'on/presentaci\'on que facilita extensi\'on y pruebas unitarias.
  \item Plataforma desplegada p\'ublicamente, accesible sin instalaci\'on.
\end{itemize}

\subsection{Limitaciones}

\begin{itemize}
  \item La propagaci\'on frontend usa Kepler puro sin perturbaciones en tiempo real; la propagaci\'on RK45 + J2 + drag est\'a implementada en el backend pero requiere Docker para activarse.
  \item El modelo atmosf\'erico usa densidad exponencial; la integraci\'on con NRLMSISE-00 est\'a disponible en el backend Python.
  \item La precesi\'on-nutaci\'on IAU 2006/2000A solo se aplica en el backend (v\'ia \code{astropy}); el frontend usa GMST simplificado.
\end{itemize}

\subsection{Trabajo Futuro}

\begin{itemize}
  \item Integraci\'on completa del backend FastAPI (Docker Compose).
  \item Exportaci\'on CZML para playback nativo en CesiumJS.
  \item An\'alisis de cobertura de constelaciones tipo Walker.
  \item Calculadora de rendezvous por ecuaciones de Clohessy-Wiltshire (LVLH).
  \item Validaci\'on formal contra GMAT y Orekit con tabla de errores completa.
\end{itemize}

% =============================================================================
\section*{Referencias}
\addcontentsline{toc}{section}{Referencias}

\begin{enumerate}[label={[\arabic*]}, leftmargin=2.5em]
  \item Vallado, D.A. \textit{Fundamentals of Astrodynamics and Applications}, 4th ed. Microcosm Press, 2013.
  \item Curtis, H.D. \textit{Orbital Mechanics for Engineering Students}, 4th ed. Elsevier, 2019.
  \item Bate, R.R., Mueller, D.D., White, J.E. \textit{Fundamentals of Astrodynamics}. Dover, 1971.
  \item Battin, R.H. \textit{An Introduction to the Mathematics and Methods of Astrodynamics}. AIAA, 1987.
  \item Wertz, J.R. (ed.) \textit{Space Mission Engineering: The New SMAD}. Microcosm, 2011.
  \item Brouwer, D. ``Solution of the Problem of Artificial Satellite Theory without Drag.'' \textit{Astronomical Journal}, 64, 1959.
  \item Dormand, J.R., Prince, P.J. ``A family of embedded Runge-Kutta formulae.'' \textit{Journal of Computational and Applied Mathematics}, 6(1), 1980.
  \item Celestrak --- Cat\'alogo TLE de sat\'elites activos. \url{https://celestrak.org}
  \item CesiumJS Documentation. \url{https://cesium.com/learn/cesiumjs/ref-doc/}
  \item Poliastro --- Python Astrodynamics Library. \url{https://docs.poliastro.space}
  \item IERS Conventions 2010. \url{https://www.iers.org/IERS/EN/Publications/TechnicalNotes/tn36.html}
\end{enumerate}

\end{document}
